
\section{TABLAS} \label{sec:tablas}

Las tablas se definen en el entorno \texttt{table}. Existen infinidad de posibilidades en cuanto a su formato: omitir o dibujar líneas horizontales y verticales, fusionar columnas y filas, alinear el contenido a la derecha, izquierda o centro, y demás opciones resumidas en \url{https://www.overleaf.com/learn/latex/tables}. Dado que la forma de construir una tabla directamente en código \LaTeX \ está lejos de ser cómoda e intuitiva, quizás lo más recomendable sea acudir a editores de tablas que generan automáticamente el código correspondiente y cuya interfaz es similar a la que puede encontrarse en Excel, como por ejemplo \url{https://www.tablesgenerator.com/}. Un ejemplo sencillo de tabla se muestra a continuación:

\vspace{25pt}

\begin{table}[H]
\centering
\begin{tabular}{|c|c|c|l|}
\hline
$\bm{n}$ & $\bm{a_n}$ & $\bm{a_{n+1}}$ & $\bm{\varphi} \ _{(= a_{n+1}/a_n)}$ \\ \hline\hline
\textbf{1} & 1 & 1 & 1 \\ \hline
\textbf{2} & 1 & 2 & 2 \\ \hline
\textbf{3} & 2 & 3 & 1,5 \\ \hline
\textbf{4} & 3 & 5 & 1,66666667 \\ \hline
\textbf{5} & 5 & 8 & 1,6 \\ \hline
\textbf{6} & 8 & 13 & 1,625 \\ \hline
\textbf{7} & 13 & 21 & 1,61538462 \\ \hline
\textbf{8} & 21 & 34 & 1,61904762 \\ \hline
\textbf{9} & 34 & 55 & 1,61764706 \\ \hline
\textbf{10} & 55 & - & - \\ \hline
\end{tabular}
\caption{Cinco primeros términos de la sucesión de Fibonacci}
\label{tabla:fibonacci5}
\end{table}

\vspace{5pt}

El título de la tabla se indica mediante el comando \textbackslash\texttt{caption} (este comando no solamente sirve para añadir un título a la tabla, sino que es esencial para que ésta aparezca en el índice de tablas), y, al igual que en el caso de los títulos de capítulos (ver apartado \ref{sec:referencias}), es muy recomendable añadir el comando \textbackslash\texttt{label} para poder referirse a la tabla en cuestión en partes posteriores (o anteriores) del documento.


\subsection{Posicionamiento}

Las tablas (al igual que otros elementos como imágenes o diagramas, como se verá más adelante) pueden posicionarse en distintos lugares de la página y en distintas posiciones con respecto al texto. La forma más común de situar una tabla es inmediatamente después de un párrafo y centrada en la página (como en el caso de la tabla \ref{tabla:fibonacci5}), lo que se consigue indicando \texttt{[H]} al iniciar el entorno \texttt{table} y añadiendo el comando \textbackslash\texttt{centering}, respectivamente. Una guía que recopila las distintas opciones en lo que se refiere al posicionamiento de tablas e imágenes puede consultarse en \url{https://www.overleaf.com/learn/latex/positioning_images_and_tables}.


\newpage
\subsection{Entorno \textbackslash\texttt{longtable}}

En el caso de que una tabla sea demasiado larga como para caber en una única página se puede utilizar el entorno \texttt{longtable}, mediante el cual \LaTeX \ secciona la tabla de forma automática en tantas partes como sea necesario.

\vspace{10pt}

\begin{longtable}[H]{|c|c|c|l|}
\hline
$\bm{n}$ & $\bm{a_n}$ & $\bm{a_{n+1}}$ & $\bm{\varphi} \ _{(= a_{n+1}/a_n)}$ \\ \hline\hline
\endhead
\textbf{1} & 1 & 1 & 1 \\ \hline
\textbf{2} & 1 & 2 & 2 \\ \hline
\textbf{3} & 2 & 3 & 1,5 \\ \hline
\textbf{4} & 3 & 5 & 1,66666667 \\ \hline
\textbf{5} & 5 & 8 & 1,6 \\ \hline
\textbf{6} & 8 & 13 & 1,625 \\ \hline
\textbf{7} & 13 & 21 & 1,61538462 \\ \hline
\textbf{8} & 21 & 34 & 1,61904762 \\ \hline
\textbf{9} & 34 & 55 & 1,61764706 \\ \hline
\textbf{10} & 55 & 89 & 1,61818182 \\ \hline
\textbf{11} & 89 & 144 & 1,61797753 \\ \hline
\textbf{12} & 144 & 233 & 1,61805556 \\ \hline
\textbf{13} & 233 & 377 & 1,61802575 \\ \hline
\textbf{14} & 377 & 610 & 1,61803714 \\ \hline
\textbf{15} & 610 & 987 & 1,61803279 \\ \hline
\textbf{16} & 987 & 1597 & 1,61803445 \\ \hline
\textbf{17} & 1597 & 2584 & 1,61803381 \\ \hline
\textbf{18} & 2584 & 4181 & 1,61803406 \\ \hline
\textbf{19} & 4181 & 6765 & 1,61803396 \\ \hline
\textbf{20} & 6765 & 10946 & 1,618034 \\ \hline
\textbf{21} & 10946 & 17711 & 1,61803399 \\ \hline
\textbf{22} & 17711 & 28657 & 1,61803399 \\ \hline
\textbf{23} & 28657 & 46368 & 1,61803399 \\ \hline
\textbf{24} & 46368 & 75025 & 1,61803399 \\ \hline
\textbf{25} & 75025 & 121393 & 1,61803399 \\ \hline
\textbf{26} & 121393 & 196418 & 1,61803399 \\ \hline
\textbf{27} & 196418 & 317811 & 1,61803399 \\ \hline
\textbf{28} & 317811 & 514229 & 1,61803399 \\ \hline
\textbf{29} & 514229 & 832040 & 1,61803399 \\ \hline
\textbf{30} & 832040 & 1346269 & 1,61803399 \\ \hline
\textbf{31} & 1346269 & 2178309 & 1,61803399 \\ \hline
\textbf{32} & 2178309 & 3524578 & 1,61803399 \\ \hline
\textbf{33} & 3524578 & 5702887 & 1,61803399 \\ \hline
\textbf{34} & 5702887 & 9227465 & 1,61803399 \\ \hline
\textbf{35} & 9227465 & 14930352 & 1,61803399 \\ \hline
\textbf{36} & 14930352 & 24157817 & 1,61803399 \\ \hline
\textbf{37} & 24157817 & 39088169 & 1,61803399 \\ \hline
\textbf{38} & 39088169 & 63245986 & 1,61803399 \\ \hline
\textbf{39} & 63245986 & 102334155 & 1,61803399 \\ \hline
\textbf{40} & 102334155 & 165580141 & 1,61803399 \\ \hline
\textbf{41} & 165580141 & 267914296 & 1,61803399 \\ \hline
\textbf{42} & 267914296 & 433494437 & 1,61803399 \\ \hline
\textbf{43} & 433494437 & 701408733 & 1,61803399 \\ \hline
\textbf{44} & 701408733 & 1134903170 & 1,61803399 \\ \hline
\textbf{45} & 1134903170 & 1836311903 & 1,61803399 \\ \hline
\textbf{46} & 1836311903 & 2971215073 & 1,61803399 \\ \hline
\textbf{47} & 2971215073 & 4807526976 & 1,61803399 \\ \hline
\textbf{48} & 4807526976 & 7778742049 & 1,61803399 \\ \hline
\textbf{49} & 7778742049 & 1,2586E+10 & 1,61803399 \\ \hline
\textbf{50} & 1,2586E+10 & - & - \\ \hline
\caption{Cincuenta primeros términos de la sucesión de Fibonacci}
\label{tabla:fibonacci50}
\end{longtable}

Existen distintas alternativas en cuanto a qué elementos incluir tanto en la primera como en la última línea de cada sección de tabla (en el caso de la tabla \ref{tabla:fibonacci50} se ha elegido repetir la primera línea en cada una de sus secciones), que pueden consultarse en \url{https://texblog.org/2011/05/15/multi-page-tables-using-longtable/}.

%%%%%%%%%%%%%%%%%%%%%%%%%%%%%%%%%%%%%%%%%%%%%%%%%%

